% 2009 年湖北卷文科第 17 题：矩形场地与旧墙。
% 来源 PDF：普通高考/2009/2009湖北文.pdf，SHA-256 f60ba7a5a2fef3449fba9474fa973c616aa52a15926b3a97f0dda6b909f4c861。
% 原页：第 1 页第 3 栏中部，书页码 525。
% 重绘依据：照原页方位，旧墙在上方并画成带竖缝的墙身，左右两侧为新墙，
%   与旧墙相对的新墙在中间留出宽度 2 m 的进出口（缺口）；x 为所利用旧墙的长度，
%   用双箭头标在矩形内侧上方。图不按比例，只表示位置关系。
\begin{tikzpicture}[x=1cm, y=1cm, semithick,
  every node/.style={inner sep=1.5pt, font=\small}]

  % 旧墙墙身
  \draw[thick] (-0.75,0.40) -- (6.75,0.40);
  \draw[thick] (-0.75,0.06) -- (6.75,0.06);
  \foreach \xx in {-0.45,-0.15,0.15,...,6.75} {
    \draw (\xx,0.06) -- (\xx,0.40);
    \draw (\xx,0.40) -- (\xx,0.52);
    \draw (\xx,0.06) -- (\xx,-0.06);
  }

  % 左右两侧新墙
  \draw (0,0.06) -- (0,-2.75);
  \draw (6,0.06) -- (6,-2.75);

  % 与旧墙相对的新墙，中部留 2 m 进出口
  \draw (0,-2.75) -- (2.65,-2.75);
  \draw (3.35,-2.75) -- (6,-2.75);

  % 旧墙长度 x
  \draw[{Stealth}-{Stealth}] (0,-0.85) -- (6,-0.85)
    node[midway, above] {$x$};
\end{tikzpicture}
